% ch6_b §6.5–6.7 (书页 181–189 / p195–p203)
\section{准经典动力学}

在任何涉及电子在晶格中运动的问题里，我们都可以着手去求解等效 Schrödinger 方程 (6.30)。但如果 $\mathcal{E}(\mathbf{k})$ 是 $\mathbf{k}$ 的一个复杂函数，这样做可能会导致相当棘手的数学。

应用对应原理（correspondence principle）可以避开其中大部分困难。众所周知，Schrödinger 方程的波包解与经典粒子遵循同样的轨迹，后者服从由相应的经典哈密顿量导出的运动方程。按照量子力学的Schrödinger 表述，波动方程中的哈密顿量\emph{算符}是这样得到的：在经典哈密顿量\emph{函数}中，用 $-i\hbar\nabla$ 代替经典动量变量 $\mathbf{p}$。把上述步骤反过来，我们可以在等效哈密顿量算符 $\mathcal{E}(-i\nabla)$ 中，把算符 $-i\nabla$ 换成动量变量 $\mathbf{p}/\hbar$，并把所得结果称为等效经典哈密顿量函数。于是，
\begin{equation*}
\mathcal{E}(-i\nabla)+\mathcal{U}(\mathbf{r})\;\rightarrow\;\mathcal{E}(\mathbf{p}/\hbar)+\mathcal{U}(\mathbf{r})=\mathcal{H}(\mathbf{r},\mathbf{p}). \tag{6.34}
\end{equation*}
我们的电子波包将表现得如同经典点粒子，服从由这个哈密顿量函数生成的运动方程。

Hamilton 方程可以符号式地写成
\begin{equation*}
\dot{\mathbf{r}}=\frac{\partial\mathcal{H}}{\partial\mathbf{p}},\quad \dot{\mathbf{p}}=-\frac{\partial\mathcal{H}}{\partial\mathbf{r}}. \tag{6.35}
\end{equation*}
其中第一式便成为一个关于速度的方程，
\begin{align*}
\mathbf{v}&=\frac{\partial\mathcal{H}}{\partial\mathbf{p}}=\frac{\partial}{\partial\mathbf{p}}\bigl\{\mathcal{E}(\mathbf{p}/\hbar)+\mathcal{U}(\mathbf{r})\bigr\}\\
&=\frac{\partial\mathcal{E}(\mathbf{p}/\hbar)}{\partial\mathbf{p}}\\
&=\frac{1}{\hbar}\frac{\partial\mathcal{E}(\mathbf{k})}{\partial\mathbf{k}}. \tag{6.36}
\end{align*}
用我们更常用的变量 $\mathbf{k}$ 表出。这正是式 (6.2)，至此终于得到证明。我们看到
\begin{equation*}
\hbar\mathbf{k}=\mathbf{p} \tag{6.37}
\end{equation*}
的确扮演着经典动量的角色；尽管它并不是一个唯一确定的变量，它却标记着波包恰好主要集中于其上的那个“态”。

Hamilton 方程中的第二式可以写成
\begin{equation*}
\hbar\dot{\mathbf{k}}=-\frac{\partial\mathcal{H}}{\partial\mathbf{r}}=-\nabla\mathcal{U}(\mathbf{r}). \tag{6.38}
\end{equation*}
但如果 $\mathcal{U}(\mathbf{r})$ 譬如说是电子在固定静电场中的势能，那么右端就简单地是作用在电子上的经典力。例如，在电场 $\mathbf{E}$ 中，我们有
\begin{equation*}
\hbar\dot{\mathbf{k}}=e\mathbf{E}. \tag{6.39}
\end{equation*}
这就证明了我们的第二条动力学定律 (6.5)。变量 $\hbar\mathbf{k}$ 之所以扮演动量的角色，在于恒力的效果是使 $\hbar\mathbf{k}$ 以恒定速率增大。这里再次看到，$\mathbf{k}$ 不唯一这一事实并无妨碍，因为进入这些方程的只是它的导数。

至于磁场 $\mathbf{H}$ 的效应，我们假定 \emph{Lorentz 力}（Lorentz force）方程成立，即
\begin{equation*}
\hbar\dot{\mathbf{k}}=e\Bigl(\mathbf{E}+\frac{1}{c}\,\mathbf{v}\wedge\mathbf{H}\Bigr), \tag{6.40}
\end{equation*}
其中 $\mathbf{v}$ 是波包的速度，由式 (6.36) 算得；$c$ 是光速。这个公式是从式 (6.39) 类比而来的，但它的正当推导要困难得多。人们相信，直到非常强的磁场它都成立得相当好（见 §9.8）。

\section{质量张量：电子与空穴}

让我们把运动学方程和动力学方程 (6.36) 与 (6.39) 结合起来，计算一个电子的\emph{加速度}，
\begin{align*}
\dot{\mathbf{v}}&=\frac{\partial}{\partial t}\biggl\{\frac{1}{\hbar}\frac{\partial\mathcal{E}(\mathbf{k})}{\partial\mathbf{k}}\biggr\}\\
&=\frac{1}{\hbar}\frac{\partial}{\partial\mathbf{k}}\biggl\{\frac{1}{\hbar}\frac{\partial\mathcal{E}(\mathbf{k})}{\partial\mathbf{k}}\biggr\}\cdot\frac{\hbar\,\partial\mathbf{k}}{\partial t}\\
&=\frac{1}{\hbar^{2}}\frac{\partial^{2}\mathcal{E}(\mathbf{k})}{\partial\mathbf{k}\,\partial\mathbf{k}}\cdot e\mathbf{E}. \tag{6.41}
\end{align*}
%FIG{102}: 质量沿"谷"的方向很大。

若把它与常规的 Newton 方程
\begin{equation*}
m\dot{\mathbf{v}}=e\mathbf{E}, \tag{6.42}
\end{equation*}
相比较，我们看到电子质量的等效物如今是一个张量的逆：
\begin{equation*}
\frac{1}{m}\;\rightarrow\;\frac{1}{\hbar^{2}}\frac{\partial^{2}\mathcal{E}(\mathbf{k})}{\partial\mathbf{k}\,\partial\mathbf{k}}. \tag{6.43}
\end{equation*}
这个张量是最接近于电子的动力学质量的等效物。例如，在半导体中，能量面可以具有 (4.35) 的形式，即
\begin{equation*}
\mathcal{E}(\mathbf{k})=\frac{\hbar^{2}k_{1}^{2}}{2m_{1}}+\frac{\hbar^{2}k_{2}^{2}}{2m_{2}}+\frac{\hbar^{2}k_{3}^{2}}{2m_{3}} \tag{6.44}
\end{equation*}
它参照某个局域主轴系写出。电场在不同方向上对电子加速的效果于是将大不相同。如果能量面沿 $k_{1}$ 方向拉长成椭球，即若 $m_{1}\gg m_{2}\sim m_{3}$，那么沿“谷”方向测得的质量就显得比横向方向大得多。在计算各种态对电子平均迁移率的贡献时，这一点很重要。施主杂质能级的简单类氢公式 (6.33) 也必须加以修正，以顾及每个谷的等效哈密顿量的各向异性，以及许多半导体导带底处同一能量诸谷的多重性。

\emph{逆质量张量}（inverse mass tensor）不一定是正定的。在（如贵金属 Cu、Ag、Au 中）费米面与区边界相接触之处，我们可以近似地把它表成 (6.44) 的形式，但取 $m_{3}$ 为\emph{负}，从而给出双曲型能量面（在 §3.3 的重复区图式中，可以让它们恰好穿过区边界）。从经典的观点看，$\mathbf{k}$ 空间中这样一点上电子的动力学行为会显得非常奇怪。

它们被来自区边界的强 Bragg 反射深刻地改变。实际上，电子倾向于沿着晶格平面传播，仿佛穿过一系列波导一般。

%FIG{103}: 双曲面状能量面（"颈"）。

%FIG{104}: 近满带带顶处的"空穴袋"。

更常见的情形是像半导体或半金属中那样的\emph{近满带}（nearly full band）。由于 $\mathcal{E}(\mathbf{k})$ 在重复区中是连续的周期函数，被占据的态将几乎一直填充到 $\mathcal{E}(\mathbf{k})$ 的某个明确的极大值。若这个极大值出现在 $\mathbf{k}_{0}$ 处，那么作为 $\mathcal{E}(\mathbf{k})$ 的梯度的 $\mathbf{v}_{\mathbf{k}}$ 在 $\mathbf{k}=\mathbf{k}_{0}$ 处为零，而且在 $\mathbf{k}_{0}$ 附近的一切方向上，$\partial^{2}\mathcal{E}/\partial\mathbf{k}\,\partial\mathbf{k}$ 的主分量都是\emph{负}的。

具有\emph{负有效质量}（negative effective mass）的电子，其性质与寻常粒子相差太远，以至于更便当的做法是把论证反过来说，去讨论满带中少数\emph{空穴}（hole）的性质。我们在 §4.6 中已经指出，这种空穴的数目可以相当简单地用 费米–狄拉克 统计算出，只需把能量\emph{向下}量度；现在我们要证明，如果把空穴当作\emph{带正电}的粒子来处理，同一个能量约定将以经典上可理解的形式描述它们的动力学性质。

%FIG{105}: (a) 电子分布中的一个空隙。(b) 相应的"空穴"，能量标度已倒置。

我们首先注意到，“处于态 $\mathbf{k}$ 的空穴”必须由一个行列式来描述，这个行列式包含带中除态 $\psi_{\mathbf{k}}$ 之外的所有波函数。我们还可以更进一步，实际构造出这种行列式的波包。但需要注意的要点是，空穴的“速度”与那个缺失态的速度相同；波包将以与空穴两侧邻居相同的速率被携带着通过空间：
\begin{equation*}
\mathbf{v}_{\mathbf{k}}=\frac{1}{\hbar}\frac{\partial\mathcal{E}(\mathbf{k})}{\partial\mathbf{k}}. \tag{6.45}
\end{equation*}
现在来考虑与一个空穴相联系的电流。满带的电流为零。处于态 $\mathbf{k}$ 的一个电子的电流是 $-|e|\mathbf{v}_{\mathbf{k}}$。因此，缺了一个电子的满带，其电流必定是
\begin{equation*}
\mathbf{J}_{h}=0-(-|e|\mathbf{v}_{\mathbf{k}})=|e|\mathbf{v}_{\mathbf{k}}. \tag{6.46}
\end{equation*}
可见空穴的行为就像一个\emph{带正电的粒子}（positively charged particle）。

同样，空穴的动力学必定与它两侧各态的动力学相同。速度将以式 (6.41) 所定义的速率改变。若用电流来表述，我们有
\begin{align*}
|e|\dot{\mathbf{v}}&=|e|\,\frac{1}{\hbar^{2}}\frac{\partial^{2}\mathcal{E}(\mathbf{k})}{\partial\mathbf{k}\,\partial\mathbf{k}}\cdot(-|e|)\mathbf{E}\\
&=|e|\biggl\{-\frac{1}{\hbar^{2}}\frac{\partial^{2}\mathcal{E}(\mathbf{k})}{\partial\mathbf{k}\,\partial\mathbf{k}}\biggr\}\cdot|e|\mathbf{E}, \tag{6.47}
\end{align*}
这就是一个电荷为正、大小为 $|e|$ 的粒子的加速度方程，只是其质量为 $m^{*}$，其中
\begin{equation*}
\frac{1}{m^{*}}\sim-\frac{1}{\hbar^{2}}\frac{\partial^{2}\mathcal{E}(\mathbf{k})}{\partial\mathbf{k}\,\partial\mathbf{k}}. \tag{6.48}
\end{equation*}
但由于我们是在带顶附近，曲率的所有分量都是负的，所以 $m^{*}$ 现在是一个正量。于是，空穴的动力学性质未必各向同性，但它们基本上就是具有正质量的经典粒子的动力学性质。

%FIG{106}: 小重叠的能带：(a) 扩展区图式；(b) 重复区图式。

价带带顶的空穴导电，当然是半导体的一个特征性质。在 Si 和 Ge 中，区中心的能量面被自旋–轨道相互作用所分裂（参看图 66），我们实际上可以观察到“轻”“重”两种空穴，它们对应于 $\mathcal{E}(\mathbf{k})$ 的两支，曲率不同而在极大值处相接触。

%FIG{107}: (a) 电子分布中的一个空球变成 (b) 一个"空穴球"。

显然，当我们离开极大值位置时，空穴的速度会增大，也就是说，随 $\mathcal{E}(\mathbf{k})$ 的减小而增大。同样方便的做法是把空穴的动能看成 $\mathcal{E}_{h}=\mathcal{E}_{v}-\mathcal{E}(\mathbf{k})$，即从价带顶这个原点向外增大。这就是 §4.6 中引入的约定。

对于每单位晶胞含偶数个电子的半金属（如 Bi），如果存在向更高能带的小重叠、有少量电子溢了出去，我们也会遇到类似的情形。这时在近满带中会留下等数目的空穴——它们将由费米面的若干碎片包住，这些碎片可以像 §3.3 中那样在重复区图式中拼合成闭合曲面。把这样一个空穴\emph{袋}（pocket）设想成少数几个带正电的粒子、其速度矢量由中心指向外，要比讨论一个电子的费米面（其速度矢量画成指向未占据空间区域的内部）来得容易。然而应当强调，把态分类为“空穴”和“电子”的做法可能失效——例如在图 103 的“颈”上。

空穴在离子晶体中能否像自由粒子那样运动，是值得怀疑的。相邻离子上价轨道之间的小重叠意味着价带非常窄（§3.4），有效质量相应地很大。这似乎允许空穴发生\emph{自陷}（self-trapping）——例如在 KCl 中，空穴定域在两个 Cl$^{-}$ 离子之间，这两个离子被拉近而形成一个稳定的 \emph{V 心}（V centre），只有靠热激发才能把空穴从中释放出来。这样的\emph{小极化子}（small polaron）（§8.3）对晶格的颗粒性如此敏感，以至于只能靠从一个晶胞到另一个晶胞的“跳跃”来移动。由于这个原因，“价带”的概念在离子晶体中没有什么意义（§4.2）。

\section{激子}

存在两种载流子，电荷相反，但其他方面的性质基本上相似——这一点处于半导体应用物理的核心。我们在 §4.6 中已经证明，如何通过掺入受主杂质把空穴的数目增加到任意想要的数量，做成 \emph{p} 型样品。

%FIG{108}: 半导体中的杂质能级。

对于施主杂质，我们在 §6.4 中已经看到：由于带电杂质对额外电子的吸引，会形成正好位于导带底之下的杂质能级。对于受主杂质有一个完全等效的理论：受主杂质带负电（因为它们必须再接纳第四个价电子以维持键结构），因此吸引的是空穴。其效果是在空穴动能零点之下不远处为空穴产生束缚态——也就是说，在价带顶之上不远处。这种\emph{受主能级}（acceptor level）的“电离能”可以完全像式 (6.33) 那样近似地算出——尽管这里又出现了同样的假设，即 $\mathcal{E}(\mathbf{k})$ 是像 (6.31) 倒置过来那样的简单二次函数，这可能又过于天真。

在一个\emph{补偿样品}（compensated sample）中，施主和受主杂质都有相当的数量，一些电子可以从施主能级“落下”去，同受主上的空穴湮灭。正如式 (4.42) 中指出的，载流子的净数目及其符号将是差 $N_{d}-N_{a}$。然而值得指出的是，“腾空”了的杂质并不是电中性的客体；施主将带正电，受主将带负电，因此它们会吸引或排斥余下的载流子。

假设我们有一块本征半导体，或者一块补偿得如此精心、以致电子和空穴的数目
由热激发公式 (4.39) 给出，而且几乎相等。这两种粒子电荷相反——所以它们彼此吸引。我们可以建立类氢的波函数，让电子和空穴绕着它们共同的质心旋转；这一对粒子的能量降低将具有与 (6.33) 相同的形式，只是 $m^{*}$ 如今应取为约化有效质量。

%FIG{109}: (a) 被电子和空穴占据的施主能级与受主能级。(b) 电子与空穴结合。(c) 多余的电子被热激发进入导带。

这意味着，我们有一个略低于导带底的电子能级，与一个恰好在价带之上的空穴能级相联系——事实上，是整整一系列这样的能级，对应于类氢态的不同量子数——确切地说，是整整一条\emph{激子态}（exciton state）的带，对应于电子–空穴对共同质心的运动动能的叠加（图 110、111(a)）。

思考激子的另一种方式是借助受激原子。原子或离子上的一个电子被提升到某个激发态，譬如说 $\phi_{e}(\mathbf{r})$。于是我们构造一个如下形式的 Bloch 函数
\begin{equation*}
\psi_{\mathbf{k}}=\sum_{l}e^{i\mathbf{k}\cdot l}\,\phi_{e}(\mathbf{r}-l) \tag{6.49}
\end{equation*}
以使这个激发能在晶体中传播（图 111(b)）。乍一看，这似乎不过是导带中的一个态。但我们据以导出这一论证的紧束缚法（§3.4）是单电子理论；它没有考虑电子–电子相互作用。激子态比 (6.49) 更复杂：这个运动的电子随身携带着价带中其他电子的一种局域极化——换句话说，它随身拖着自己的空穴。因此，要把这对电子–空穴分开成独立的载流子，还需要额外的一小份能量。

%FIG{110}: 激子态。

%FIG{111}: (a) Wannier 激子。(b) Frenkel 激子。

虽然激子的这两种描述实质上是等价的，但它们各自的实际适用范围不同。在半导体中，载流子很轻、介电常数很高，\emph{Wannier 激子}（Wannier exciton）的那种扩展的类氢轨道是较好的近似；在分子晶体（如固态氩或蒽）中，从\emph{Frenkel 激子}（Frenkel exciton）模型出发更为准确，这种激子的波函数大部分位于单个晶胞之内。在 NaCl 这类离子晶体中，激子的大小居中，用“一个电子从 Cl$^{-}$ 离子转移到邻近的 Na$^{+}$ 格点”来描述就相当不错。

在半导体和绝缘体的光学光谱中可以探测到激子谱线（§8.5）；但由于这种激发是电中性的，它对电导没有直接贡献。激子虽然是可移动的，并且原则上应当能够把能量输运着穿过晶体，但要在实验上演示这一现象并不容易。
